%% Kullanma_Kilavuzu.tex — Kitapçevir Türkçe Kullanma Kılavuzu
%% Derleme (docs/ klasöründe):  xelatex Kullanma_Kilavuzu.tex   (iki kez, içindekiler için)
%% Font kurulumu latex_pipeline/templates/turkish_book.tex ile aynı desendedir.

\documentclass[11pt,a4paper]{article}

%% --- Dil & kodlama (XeLaTeX, UTF-8) ---------------------------------------
\usepackage{polyglossia}
\setmainlanguage{turkish}
\setotherlanguage{english}

%% --- Fontlar: dosya adıyla yükle (macOS TeX Live isim-lookup baypası) ----
\usepackage{fontspec}
\setmainfont{lmroman10-regular.otf}[
  BoldFont=lmroman10-bold.otf,
  ItalicFont=lmroman10-italic.otf,
  BoldItalicFont=lmroman10-bolditalic.otf]
\setsansfont{lmsans10-regular.otf}[
  BoldFont=lmsans10-bold.otf,
  ItalicFont=lmsans10-oblique.otf]
\setmonofont{lmmono10-regular.otf}[
  BoldFont=lmmonolt10-bold.otf,
  ItalicFont=lmmono10-italic.otf]

%% --- Sayfa düzeni & renkler (arayüzün indigo vurgusu) --------------------
\usepackage[margin=2.4cm]{geometry}
\usepackage{xcolor}
\definecolor{vurgu}{HTML}{4F46E5}
\definecolor{koyu}{HTML}{1F2937}
\definecolor{acikgri}{HTML}{F3F4F6}

\usepackage{booktabs}
\usepackage{longtable}
\usepackage{array}
\usepackage{enumitem}
\setlist{itemsep=2pt,topsep=3pt}
\setlength{\emergencystretch}{3em}  % bölünemeyen daktilo yazısı taşmasın
\newcolumntype{L}[1]{>{\raggedright\arraybackslash}p{#1}}

%% --- Kod/komut kutuları ---------------------------------------------------
%% Türkçe polyglossia '=' ve ':' önüne boşluk ekler (XeTeX interchar); komut
%% kutusunda bu kapatılır, komutlar kopyalanınca aynen çalışsın.
\usepackage{tcolorbox}
\tcbuselibrary{skins,breakable}
\newtcolorbox{komut}{
  colback=koyu, colframe=koyu, coltext=white,
  fontupper={\ttfamily\small\XeTeXinterchartokenstate=0\relax}, boxrule=0pt, arc=3pt,
  left=8pt,right=8pt,top=6pt,bottom=6pt, breakable}
\newtcolorbox{notkutu}[1][]{
  colback=acikgri, colframe=vurgu, coltext=koyu,
  boxrule=0pt, leftrule=3pt, arc=1pt,
  left=8pt,right=8pt,top=6pt,bottom=6pt, breakable, #1}

%% --- Başlık biçimleri -----------------------------------------------------
\usepackage{titlesec}
\titleformat{\section}{\Large\bfseries\color{vurgu}}{\thesection.}{0.5em}{}
\titleformat{\subsection}{\large\bfseries\color{koyu}}{\thesubsection}{0.5em}{}
\titlespacing{\section}{0pt}{16pt}{6pt}

\usepackage[hidelinks]{hyperref}
\hypersetup{colorlinks=true, linkcolor=vurgu, urlcolor=vurgu,
  pdftitle={Kitapçevir Kullanma Kılavuzu}, pdflang={tr}}

%% Arayüzdeki düğme ve etiketleri metin içinde belirginleştirir.
\newcommand{\dugme}[1]{\textsf{\bfseries #1}}
\newcommand{\dosya}[1]{\texttt{#1}}
%% Menü yolu oku (Kitapçevir → Bileşenler…).
\newcommand{\ok}{\(\rightarrow\)}

\begin{document}

%% ===================== KAPAK ============================================
\begin{titlepage}
\centering
\vspace*{3.5cm}
{\Huge\bfseries\color{vurgu} Kitapçevir\par}
\vspace{0.6cm}
{\Large Akademik Kitap Çeviri ve Dizgi Sistemi\par}
\vspace{0.3cm}
{\large PDF/EPUB \(\rightarrow\) Türkçe PDF\par}
\vspace{2cm}
{\huge\bfseries Kullanma Kılavuzu\par}
\vspace{1.5cm}
{\large Formülleri, şekilleri ve çapraz referansları koruyarak çevirir.\par}
\vfill
{\small Bu kılavuz kurulumu, arayüzü ve iki çeviri yolunu\\
(\emph{PDF'den TeX çıkar} ve \emph{Yerel model ile çevir}) adım adım anlatır.\par}
\vspace{0.5cm}
{\small Durum: bakım modu --- yalnız hata düzeltmeleri.\par}
\vspace{0.3cm}
{\small Ekim 2026\par}
\vspace{1cm}
\end{titlepage}

\tableofcontents
\newpage

%% ===================== GİRİŞ ============================================
\section{Kitapçevir Nedir?}

Kitapçevir, akademik kitapları (fizik, mühendislik, makine öğrenmesi vb.)
Türkçeye çevirmek ve yeniden dizmek için bir araçtır. Kitabı LaTeX'e
dönüştürür; matematiksel formüller, denklem numaraları, şekiller, kaynakça ve
çapraz referanslar çeviri boyunca korunur, yalnızca doğal dildeki metin
çevrilir. Sonuç, XeLaTeX ile dizilmiş Türkçe bir PDF'dir (EPUB kaynaklı kitapta
Türkçe EPUB de üretilir).

\begin{notkutu}
\textbf{Kilit fikir:} Metin çevrilir, \emph{teknik iskelet korunur}. Kaynak
kitapta bir denklem, bir şekil, bir referans neredeyse Türkçe çıktıda da tam
olarak oradadır. Çeviri kayıpsız hedeflenir: hiçbir cümle, dipnot ya da tablo
hücresi bilerek atlanmaz.
\end{notkutu}

\begin{notkutu}
\textbf{Platform:} Kitapçevir yalnızca \textbf{macOS çalışan, Apple Silicon
(M serisi) işlemcili bir Mac'te} çalışır; uygulama macOS 14 ya da sonrasını
ister. Arayüz başka bir sistemde (Intel Mac, Linux, Windows) açılmaz.
\end{notkutu}

\begin{notkutu}
\textbf{Durum:} Bakım modu --- yalnız hata düzeltmeleri.
\end{notkutu}

%% ===================== İKİ YOL ==========================================
\section{İki Çeviri Yolu}\label{sec:ikiyol}

Arayüzde yan yana iki kart bulunur. İkisi de aynı kaynak dilleri Türkçeye
çevirir; hangisini seçeceğiniz kalite beklentinize ve bilgisayarınızın
donanımına bağlıdır. Hedef dil her zaman Türkçedir.

\vspace{6pt}
\renewcommand{\arraystretch}{1.4}
\noindent\begin{tabular}{@{}L{3.0cm}L{6.1cm}L{5.8cm}@{}}
\toprule
 & \textbf{\color{vurgu}PDF'den TeX çıkar} (önerilen) & \textbf{Yerel model ile çevir} (Beta) \\
\midrule
\textbf{Çeviriyi kim yapar} & Siz ya da yapay zekâ asistanınız (ör.\ Claude Code) & Bu Mac'teki yerel çeviri modeli; önerilen model HY-MT 1.5 7B (MLX). MLX klasörü, Ollama, LM Studio ya da GGUF dosyasındaki bir model de seçilebilir (beta, daha az sınandı) \\
\textbf{Kaynak dil} & İngilizce, Almanca, Rusça, Fransızca, İspanyolca, İtalyanca & İngilizce, Almanca, Rusça, Fransızca, İspanyolca, İtalyanca (İngilizceden başka kaynak dillerde çeviri kalitesi seçilen modele bağlıdır; HY-MT 1.5 7B ile bu dillerden yalnız kısa Almanca ve Rusça deneme belgeleri sınandı) \\
\textbf{Çeviri modeli} & Gerekmez & Gerekir (HY-MT 1.5 7B diskte yaklaşık 15 GB) \\
\textbf{Donanım} & Her Apple Silicon Mac & Güçlü donanım ister (Bölüm~\ref{sec:donanim}) \\
\textbf{Kalite} & Çeviriyi yapana bağlı; en iyi sonuç bu yoldadır & Ön okuma / taslak kalitesi \\
\textbf{Akış} & Uygulama TeX'i çıkarır, siz çevirirsiniz, uygulama PDF'yi dizer & Uygulama baştan sona kendisi çevirir ve dizer \\
\bottomrule
\end{tabular}
\vspace{8pt}

\begin{notkutu}
\textbf{Önerilen: PDF'den TeX çıkar.} Çeviri modeli gerektirmez ve en iyi
sonucu verir. Çeviriyi bir yapay zekâ asistanına yaptırmak için arayüz, TeX
çıkarıldıktan sonra Claude Code'a verilecek hazır istemi gösterir
(bkz.\ Bölüm~\ref{sec:claude}).
\end{notkutu}

\begin{notkutu}
\textbf{Yerel model hakkında dürüst not:} ``Yerel model ile çevir'' beta
aşamasındadır: çeviri modelin gücüne bağlıdır ve çıktı gözden geçirilmemiş bir
taslaktır. Eylül 2026'daki karşılaştırmalı denemede bu yolun çevirisi, TeX
yolunda yapay zekâ asistanıyla yapılan çeviriden belirgin biçimde düşük
kalitede çıktı. Bu yol gözetimsiz ve çevrim dışı çalışır; ön okuma ve taslak
için uygundur. Yerel model kullanacaksanız en çok sınanan model olan
\textbf{HY-MT 1.5 7B}'yi seçmenizi öneririz.
\end{notkutu}

Bu yolun PDF ve EPUB çıktısının başında kısa bir \textbf{makine çevirisi
künyesi} bulunur: metnin bu bilgisayardaki hangi yerel modelle çevrildiğini,
gözden geçirilmemiş bir taslak olduğunu, kaynak dili ve çeviri tarihini yazar.
``PDF'den TeX çıkar'' çıktısında künye yoktur.

Her iki yol da aynı iskeleti paylaşır: çıkarım (PDF için Marker, EPUB için
EPUB ayrıştırıcı), Markdown\(\rightarrow\)LaTeX dönüşümü, sanitizer, kalite
denetimi (QC), AutoFix ve XeLaTeX dizgisi ortaktır. Tek fark, çeviri adımını
kimin yaptığıdır. Kitabın ilk sayfası tam sayfa görselli bir kapaksa çıktıya
kapak olarak konur; metin sayfası ya da boş sayfa kapak sayılmaz.

%% ===================== KURULUM ==========================================
\section{Kurulum}

İki kurulum yolu vardır: \textbf{uygulama (DMG) ile} (Bölüm~\ref{sec:dmg}--\ref{sec:veri};
terminal gerekmez) ve \textbf{kaynak koddan} (Bölüm~\ref{sec:kaynak};
geliştiriciler ve komut satırı kullanımı için).

\subsection{Gereksinimler}
\begin{itemize}
  \item \textbf{Bilgisayar:} Apple Silicon (M serisi) işlemcili, macOS 14 ya
        da sonrası çalışan bir Mac.
  \item \textbf{İnternet:} Kurulum sırasında bileşenler indirilir.
        Uygulamanın TeX çıkarma, yerel modelle çeviri ve dizgi adımları
        çevrim dışı çalışır (Hafif TeX yeni bir LaTeX paketine ilk kez gerek
        duyduğunda onu indirir).
  \item \textbf{Disk:} PDF desteği (Marker, torch ve Marker modelleri) birkaç
        GB yer tutar; Marker modelleri tek başına yaklaşık 3,5 GB'tır. Yerel
        model kullanılacaksa model için ayrıca yer gerekir (HY-MT 1.5 7B
        yaklaşık 15 GB; başka bir model seçilirse o modelin boyutu). Kurulum
        sihirbazı seçilen bileşenlerin boyutunu gösterir, boş diski ölçer ve
        yer yetmiyorsa kurulumu başlatmaz.
  \item \textbf{Bellek (yalnızca yerel model için):} Model çeviri sırasında
        bütünüyle belleğe yüklenir; uygulama bu Mac'in belleğini ölçüp modelin
        sığıp sığmadığına kendisi bakar (Bölüm~\ref{sec:donanim}). Ollama ya da
        LM Studio'da yüklü modeller de aynı belleği kullanır.
  \item \textbf{Kaynak koddan kurulumda ayrıca:} arm64 Python 3.12
        (Rosetta altında çalışan x86\_64 Python kabul edilmez) ve XeLaTeX
        içeren bir TeX dağıtımı (MacTeX).
\end{itemize}

\subsection{Donanım gereksinimi}\label{sec:donanim}
``PDF'den TeX çıkar'' her Apple Silicon Mac'te çalışır. ``Yerel model ile
çevir'' \textbf{güçlü donanım ister}: önerilen HY-MT 1.5 7B ile 20 sayfalık bir
kitabın çevirisinde ölçülen tepe bellek kullanımı yaklaşık 26 GB'tı (36 GB
belleği olan bir Mac'te, 29 Eylül 2026).

Uygulama bu Mac'in belleğini \textbf{ölçer}: MLX'in bildirdiği, Metal'in
önerdiği çalışma kümesi (\dosya{max\_recommended\_working\_set\_size};
belleğin kabaca üçte ikisi ile dörtte üçü) ``kullanılabilir bellek''tir.
MLX bu değeri veremezse toplam bellekten kaynaklı bir oranla hesaplanır ve bu
durum kartta yazar; hiç ölçülemezse kart açık kalır ve ``bellek ölçülemedi''
uyarısı çıkar. Sol kenardaki \dugme{Sistem} satırı çipi ve kullanılabilir
belleği gösterir.

Bir modelin tahmini gereksinimi üç kalemin toplamıdır: modelin diskte (ya da
sunucuda) ölçülen \textbf{ağırlığı}, modelin kendi üst verisinden (katman,
KV başlığı, başlık boyutu) kesin formülle hesaplanan \textbf{KV önbelleği}
(düşünerek çeviren modelde parça başına düşünme payı da eklenir) ve ağırlığın
\%25'i kadar \textbf{çalışma payı} (ölçülmemiş, kaynaklı tek varsayım).

\begin{itemize}
  \item \textbf{Genel kapı:} Bu Mac 4 bit nicemlenmiş bir 7B modeli
        (referans model) toplu iş 1 ile çalıştıramıyorsa ``Yerel model ile
        çevir'' kartı kapalıdır ve neden sayılarıyla yazar: ``Bu Mac yerel
        model için yetersiz: en az \ldots{} GB kullanılabilir bellek gerekir,
        bu Mac'te \ldots{} GB.'' Formüle göre bu 8 GB'lık Mac'lerde olur.
        ``PDF'den TeX çıkar'' kartı normal çalışır. Kurulum sihirbazında da
        HY-MT ve GGUF seçenekleri bu durumda kapalıdır.
  \item \textbf{Model başına kapı:} Toplu iş 1 ile bile sığmayan model
        listede ``bu Mac'e sığmaz (gerekli \textasciitilde\ldots{} GB,
        kullanılabilir \ldots{} GB)'' diye görünür; seçilirse \dugme{Çevir}
        soluk kalır. Formüle göre HY-MT 1.5 7B 16 GB'lık Mac'e sığmaz.
  \item \textbf{Uyarlanır toplu iş:} Bellek darsa aynı anda çevrilen parça
        sayısı sığacak kadar küçültülür; bellek yetiyorsa hiçbir ayar
        değişmez.
\end{itemize}
Komut satırı aynı hükmü verir: kart kapalıysa ya da model sığmıyorsa
\dosya{translate\_book.py} çıkarımdan önce aynı iletiyle \textbf{çıkış kodu
5} verir. Ölçüm, formül ve her varsayımın kaynağı:
\dosya{docs/DONANIM\_VARSAYIMLARI.md}.

\subsection{Uygulamayı indirin ve taşıyın}\label{sec:dmg}
\begin{enumerate}
  \item Deponun sürüm sayfasından
        (\url{https://github.com/berketez/kitap_cevir/releases/latest})
        uygulamanın DMG dosyasını (\dosya{Kitapcevir-\allowbreak<sürüm>.dmg}) indirin.
  \item DMG'yi açın ve \textbf{Kitapçevir} simgesini \textbf{Uygulamalar}
        (Applications) kısayolunun üzerine sürükleyin.
  \item Uygulamayı \textbf{Uygulamalar} klasöründen açın. DMG'nin içinden ya
        da taşımadan açarsanız açılış ekranı ``Kitapçevir'i Uygulamalar
        klasörüne taşıyıp oradan açman önerilir.'' der.
\end{enumerate}

Uygulama kendi verisini ve bileşenlerini
\dosya{\textasciitilde/Library/\allowbreak Application Support/\allowbreak Kitapçevir/}
altında tutar (Bölüm~\ref{sec:veri}); günlükler
\dosya{\textasciitilde/Library/\allowbreak Logs/\allowbreak Kitapçevir/}
altındadır. Uygulama paketinin içine hiçbir şey yazılmaz.

\subsection{İlk açılış: imzasız uygulamayı açmak}\label{sec:ilkacilis}
Kitapçevir imzasız dağıtılır: Apple geliştirici kimliğiyle imzalanmamış ve
Apple'a onaylatılmamıştır. Bu yüzden macOS ilk açılışta uygulamayı engeller.
Onay bir kez verilir; sonraki açılışlar normaldir.

\begin{itemize}
  \item \textbf{Finder'da sağ tık \ok{} Aç:} Uygulamalar klasöründe
        Kitapçevir'e sağ tıklayın (ya da Control tuşuyla tıklayın),
        \dugme{Aç}'ı seçin ve açılan pencerede yeniden \dugme{Aç}'a basın. Bu
        yol macOS 14'te çalışır.
  \item \textbf{Gerekirse Sistem Ayarları \ok{} Gizlilik ve Güvenlik \ok{}
        Yine de Aç:} macOS 15 ve sonrasında sağ tık yolu kaldırıldı.
        Uygulamayı açmayı bir kez deneyin; engellendikten sonra
        \textbf{Sistem Ayarları \ok{} Gizlilik ve Güvenlik} bölümünde,
        Güvenlik başlığı altında beliren \dugme{Yine de Aç} düğmesine basın ve
        giriş parolanızı girin. Düğme açma denemesinden sonra yaklaşık bir saat
        görünür.
\end{itemize}

\subsection{Kurulum sihirbazı}
Sihirbaz uygulamanın içindedir. Ekrandaki her sayı (boyut, hız, kalan süre)
gerçek ölçümden gelir; bilinmeyen değer ``---'' olarak gösterilir. Ekranlar
sırayla:

\begin{enumerate}
  \item \textbf{Kurulum (hoş geldin):} Bu Mac'in çipi, macOS sürümü,
        kullanılabilir belleği ve boş diski ile internet bağlantısı.
        Bileşenler yayıncılarının kendi sunucularından (ör.\ GitHub, PyPI,
        Hugging Face, datalab.to) doğrudan iner; Kitapçevir bu dosyaları
        barındırmaz.
  \item \textbf{Lisanslar:} Kitapçevir'in kendi lisansı (GNU GPL 3.0) ve PDF
        desteğinin yapay zekâ modelleri için Datalab model lisansının resmî
        metni ve bağlantısı. Datalab lisansı kabul edilmeden kurulum
        sürdürülemez: PDF desteği zorunlu bileşendir ve bu modeller olmadan
        çalışmaz (bileşenlerin lisansları: Bölüm~\ref{sec:lisans}).
  \item \textbf{Bileşenler:} Ne kurulacağını seçersiniz (aşağıdaki tablo).
        Ollama ya da LM Studio kuruluysa Kitapçevir onları kendisi bulur;
        burada kurulmaz.
  \item \textbf{Özet:} Kurulacaklar, toplam indirme, diskte kaplayacağı yer,
        boş disk ve kurulum yeri. Boş disk yetmiyorsa ne kadar daha gerektiği
        yazar.
  \item \textbf{Kurulum:} Genel ilerleme yalnızca indirmede, gerçek baytlarla
        gösterilir (hız ve tahmini kalan süreyle); kurma adımlarında çubuk
        yoktur, geçen süre yazar. \dugme{Duraklat}, \dugme{İptal} ve
        \dugme{Sürdür} düğmeleri vardır: yarım inen dosyalar saklanır, kurulum
        kaldığı yerden sürer.
  \item \textbf{Kurulum tamamlandı:} \dugme{Kitapçevir'i aç} düğmesiyle
        arayüz aynı pencerede açılır.
\end{enumerate}

\vspace{4pt}
\renewcommand{\arraystretch}{1.3}
\noindent\begin{tabular}{@{}L{3.3cm}L{3.5cm}L{8.2cm}@{}}
\toprule
\textbf{Bölüm} & \textbf{Bileşen} & \textbf{Not} \\
\midrule
Zorunlu (kilitli) & Çekirdek & Python, arayüz, EPUB yolu, derleme \\
 & MLX & Apple Silicon hesaplama kitaplığı \\
 & PDF desteği & Marker + torch + Marker modelleri (Datalab lisansı) \\
\midrule
TeX dağıtımı (birini seçin) & Hafif TeX (tectonic) & Varsayılan. Paketleri gerektikçe indirir; kurulumda şablon bir kez derlenerek ısıtılır \\
 & Tam TeX (MacTeX) & Büyük indirme; macOS Yükleyici açılır ve yönetici parolası ister \\
 & Kurulu olanı kullan & Yalnızca bu Mac'te XeLaTeX içeren bir TeX bulunursa görünür; seçilmeden önce şablon onunla derlenerek sınanır \\
\midrule
İsteğe bağlı & HY-MT 1.5 7B & Önerilen yerel çeviri modeli (yaklaşık 15 GB). Tencent lisansının kabulü gerekir; bu Mac'e sığmıyorsa seçilemez ve nedeni yazar \\
 & GGUF desteği & llama-cpp-python (Metal); GGUF dosyalarını doğrudan açmak için \\
\bottomrule
\end{tabular}
\vspace{6pt}

Kurulum sürerken uygulamayı kapatırsanız indirilen kısım saklanır ve bir
sonraki açılışta kurulum kaldığı yerden sürer.

\subsection{Bileşenler ekranı, güncelleme ve kaldırma}\label{sec:bilesenler}
\textbf{Kitapçevir \ok{} Bileşenler\ldots} menüsü (Komut + virgül) ya da
arayüzün sol kenarındaki \dugme{Bileşenler} bağlantısı Bileşenler ekranını
açar. Tabloda her bileşenin durumu, sürümü ve diskte kapladığı yer
(gerçek ölçüm) yazar.

\begin{itemize}
  \item \dugme{Onar}: bileşenin dosyalarını SHA-256 ile doğrular, bozuk ya da
        eksik olanları yeniden indirir ve yeniden kurar.
  \item \dugme{İndir} / \dugme{Kaldır}: HY-MT 1.5 7B ve GGUF desteği sonradan
        indirilebilir ya da kaldırılabilir. Kaldırırken kitaplarınız ve çeviri
        belleğiniz silinmez.
  \item \dugme{Türünü değiştir\ldots}: Hafif TeX, Tam TeX ve kurulu TeX
        arasında geçiş.
  \item \dugme{Denetle}: yeni sürüm yalnızca bu düğmeyle denetlenir,
        kendiliğinden denetlenmez ve kurulmaz. Yeni sürüm varsa
        \dugme{Sürüm sayfasını aç} sürüm sayfasını açar; yeni DMG'deki
        uygulamayı eskisinin üzerine sürüklersiniz. Kitaplarınız ve
        ayarlarınız uygulamadan ayrı tutulduğu için güncellemeden etkilenmez.
  \item \dugme{Kaynak kodu}, \dugme{Lisanslar} ve \dugme{Günlükleri göster}
        bağlantıları.
  \item \dugme{Kitapçevir'i kaldır\ldots}: silinecek klasörler boyutlarıyla
        listelenir. Varsayılan olarak \dosya{veri/} klasörü \textbf{korunur}:
        kitaplarınız, çeviri belleğiniz ve modeller (sihirbazın kurduğu
        HY-MT 1.5 7B dahil; panel adını ve diskte kapladığı gerçek yeri
        ayrıca gösterir). Yalnız HY-MT'yi silmek için ``Şunu da sil: HY-MT
        \ldots'' kutusunu, verinin tamamını silmek için ``Kitaplarımı,
        belleği ve modeller klasörünün tamamını da sil'' kutusunu
        işaretleyin. Kaldırma bitince uygulamayı kapatıp (Komut + Q) Uygulamalar
        klasöründen Çöp Sepeti'ne sürüklersiniz. MacTeX sistem geneline
        kurulduğu için kaldırılmaz; kaldırmak için TeX Users Group'un kaldırma
        yönergesine bakın.
\end{itemize}

Bir iş sürerken bileşen değiştirilemez: önce işin bitmesini bekleyin ya da
işi iptal edin.

\textbf{Elle kaldırma} (uygulama açılamıyorsa): şu klasörleri silin, sonra
uygulamayı Çöp Sepeti'ne sürükleyin. Kitaplarınızı saklamak istiyorsanız önce
\dosya{veri/} alt klasörünü başka bir yere taşıyın.
\begin{itemize}
  \item \dosya{\textasciitilde/Library/\allowbreak Application Support/\allowbreak Kitapçevir/}
  \item \dosya{\textasciitilde/Library/\allowbreak Logs/\allowbreak Kitapçevir/}
  \item varsa
        \dosya{\textasciitilde/Library/\allowbreak WebKit/\allowbreak io.github.\allowbreak berketez.\allowbreak kitapcevir},
        \dosya{\textasciitilde/Library/\allowbreak Caches/\allowbreak io.github.\allowbreak berketez.\allowbreak kitapcevir},
        \dosya{\textasciitilde/Library/\allowbreak HTTPStorages/\allowbreak io.github.\allowbreak berketez.\allowbreak kitapcevir},
        \dosya{\textasciitilde/Library/\allowbreak Preferences/\allowbreak io.github.\allowbreak berketez.\allowbreak kitapcevir.plist} ve
        \dosya{\textasciitilde/Library/\allowbreak Saved Application State/\allowbreak io.github.\allowbreak berketez.\allowbreak kitapcevir.savedState}
\end{itemize}

\subsection{Veri klasörü: kitaplar, modeller, çeviri belleği}\label{sec:veri}
DMG kurulumunda kullanıcı verisi uygulamanın kodundan ayrı, tek bir kökte
durur:
\dosya{\textasciitilde/Library/\allowbreak Application Support/\allowbreak Kitapçevir/\allowbreak veri/}.
Kaynak koddan kurulumda aynı klasörler proje klasörünün içindedir.

\vspace{4pt}
\renewcommand{\arraystretch}{1.3}
\noindent\begin{tabular}{@{}L{5.6cm}L{4.6cm}L{4.8cm}@{}}
\toprule
\textbf{İçerik} & \textbf{DMG kurulumu} (\dosya{Kitapçevir/} altında) & \textbf{Kaynak koddan} (proje klasöründe) \\
\midrule
Kitap klasörleri & \dosya{veri/output/\allowbreak<kitap\_adı>/} & \dosya{output/\allowbreak<kitap\_adı>/} \\
Kendi eklediğiniz modeller (MLX klasörleri, GGUF dosyaları); HY-MT de buraya iner & \dosya{veri/models/} & \dosya{models/} \\
Çeviri belleği & \dosya{veri/data/\allowbreak translation\_memory.db} & \dosya{data/\allowbreak translation\_memory.db} \\
İsteğe bağlı ayar dosyaları & \dosya{veri/config/} & \dosya{config/} \\
\bottomrule
\end{tabular}
\vspace{6pt}

Arayüzdeki \dugme{Klasörü aç} düğmesi kitabın klasörünü doğrudan açar. Bu
kılavuzda \dosya{models/} ve \dosya{config/} dendiğinde DMG kurulumunda
\dosya{veri/models/} ve \dosya{veri/config/} anlaşılmalıdır.

\subsection{Kaynak koddan kurulum (geliştiriciler için)}\label{sec:kaynak}
\begin{komut}
git clone https://github.com/berketez/kitap\_cevir.git kitapcevir\\
cd kitapcevir\\
python3.12 -m venv venv\\
source venv/bin/activate\\
pip install --no-deps -r requirements.txt\\
brew install --cask mactex \ \ \# XeLaTeX yoksa
\end{komut}

Sanal ortamın adı \dosya{venv} olmalıdır: arayüzü başlatma komutu bu adı
kullanır. PDF çıkarım motoru Marker, ilk PDF'de kendi modellerini (yaklaşık
3,5 GB) kendiliğinden indirir; bu yüzden ilk PDF her zamankinden uzun sürer.
Bu modellerin ağırlıkları Datalab'ın model lisansıyla sunulur; metin:
\url{https://github.com/datalab-to/marker} (\dosya{MODEL\_LICENSE}).

\textbf{Yalnızca ``Yerel model ile çevir'' için} önerilen modeli indirip MLX
biçimine dönüştürün (yaklaşık 15 GB):

\begin{komut}
./venv/bin/mlx\_lm.convert --hf-path tencent/HY-MT1.5-7B \textbackslash\\
\ \ --mlx-path models/hymt-7b-mlx
\end{komut}

\begin{notkutu}
\textbf{Yalnızca TeX yolunu kullanacaksanız bu adımı atlayın.} TeX yolu
çeviri modeli gerektirmez; PDF çıkarımı (Marker) ve XeLaTeX yeterlidir.
Başka bir model kullanacaksanız (MLX klasörü, Ollama, LM Studio ya da GGUF
dosyası) HY-MT'yi indirmek gerekmez; hazırlık Bölüm~\ref{sec:modeller}'de
anlatılır.
\end{notkutu}

%% ===================== ARAYÜZ ===========================================
\section{Uygulamayı Başlatma}

\textbf{DMG kurulumunda} Uygulamalar klasöründen \textbf{Kitapçevir}'i açın.
Arayüz uygulamanın kendi penceresinde açılır; \textbf{Görünüm \ok{} Tarayıcıda
Aç} aynı sayfayı varsayılan tarayıcıda açar. Uygulamayı Komut + Q ile kapatın;
bir iş sürerken kapatırsanız onay istenir ve iş iptal edilir.

\textbf{Kaynak koddan} proje klasöründe şu komutu çalıştırın:

\begin{komut}
./venv/bin/streamlit run app.py
\end{komut}

Arayüz tarayıcıda \dosya{http://127.0.0.1:8501} adresinde açılır (tarayıcı
kendiliğinden açılmazsa terminalde yazan adresi kullanın). Kapatmak için
terminalde \texttt{Ctrl+C}'ye basın. İki kurulumda da arayüz yalnızca bu
bilgisayardan erişilebilir; en büyük yükleme boyutu 500 MB'tır.

\subsection{Arayüzün bölümleri}

\renewcommand{\arraystretch}{1.35}
\noindent\begin{tabular}{@{}L{3.2cm}L{12.2cm}@{}}
\toprule
\textbf{Bölge} & \textbf{İçerik} \\
\midrule
\textbf{Sol kenar} & \dugme{Kitap adı} (isteğe bağlı): çıktı klasörünün adı.
Boş bırakılırsa dosya adından çıkarılır; aynı adda başka bir kitap varsa
ayırmak için kullanılır. \dugme{Sistem}: ``macOS · Apple Silicon'', çip ve
kullanılabilir bellek, seçili çeviri modelinin hazır olup olmadığı. DMG
kurulumunda \dugme{Bileşenler} bağlantısı (Bölüm~\ref{sec:bilesenler}).
\dugme{Hakkında}: proje ve Tencent'le bağlantısızlık beyanı,
\dugme{Kaynak kodu} bağlantısı. \\
\textbf{Orta, ``Çeviri''} & Kitap yükleme alanı (PDF ya da EPUB). Kitap
yüklenince dosya kartı (biçim, boyut, PDF'de sayfa sayısı) ve iki çeviri
kartı görünür: \dugme{PDF'den TeX çıkar} ve \dugme{Yerel model ile çevir}
(Beta). İş sürerken ilerleme paneli, iş bitince sonuç burada görünür. \\
\textbf{Sağ, ``Araçlar''} & \dugme{Terimler}, \dugme{Bellek} ve
\dugme{Karşılaştırma} sekmeleri (bkz.\ Bölüm~\ref{sec:araclar}). \\
\bottomrule
\end{tabular}
\vspace{6pt}

Aynı anda tek iş çalışır; bir iş sürerken çeviri düğmeleri, kitap adı ve
kitap yükleme kapalıdır. Sayfayı yenileseniz de çalışan iş sürer ve ilerlemesi
görünmeye devam eder.

%% ===================== TEX YOLU =========================================
\section{PDF'den TeX Çıkar --- Adım Adım (Önerilen)}

Bu yol üç adımdır: \textbf{çıkar} \(\rightarrow\) \textbf{çevir}
\(\rightarrow\) \textbf{derle}. Çeviri modeli gerekmez.

\subsection{Adım 1 --- TeX'i çıkarın}
\begin{enumerate}
  \item Orta alana kitabı (PDF ya da EPUB) yükleyin. İsterseniz önce sol
        kenarda kitap adını yazın.
  \item \dugme{PDF'den TeX çıkar} kartında \dugme{1 · TeX'i çıkar} düğmesine
        basın.
  \item İş bitince ``TeX çıkarıldı'' iletisi ve \dugme{EXPORT.tex'i indir}
        düğmesi çıkar. Dosya kitap klasöründe \dosya{EXPORT.tex} olarak
        kaydedilmiştir (Bölüm~\ref{sec:veri}).
  \item Kartta beliren \dugme{Klasörü aç} düğmesi kitap klasörünü Finder'da
        açar. Kart ayrıca çevrilecek dosyanın ve bu kitaba özel çeviri
        talimatının tam yolunu gösterir.
\end{enumerate}

\begin{notkutu}
\textbf{\dosya{EXPORT.tex}'in başındaki \dosya{\% KITAPCEVIR-META} satırlarını
silmeyin.} Derleme adımı kitabı bu satırlarla tanır; dosyayı başka bir
klasöre taşısanız bile kimlik bu sayede kaybolmaz.
\end{notkutu}

\subsection{Adım 2 --- EXPORT.tex'i çevirin}
Çeviriyi ister siz yapın ister bir yapay zekâ asistanı. Çevrilmiş dosyayı
\dosya{EXPORT\_TR.tex} adıyla \textbf{aynı klasöre} kaydedin. Çevirirken şu
kurallara uyulur:

\begin{itemize}
  \item Yalnızca \textbf{doğal dildeki metin ve başlıklar} çevrilir.
  \item \textbf{Dokunulmaz:} matematik (\dosya{\$...\$},
        \dosya{\textbackslash begin\{equation\}}),
        \dosya{\textbackslash ref}/\dosya{\textbackslash cite}/\dosya{\textbackslash label}
        argümanları, şekiller, dosya yolları, ortam adları ve
        \dosya{\% KITAPCEVIR-META} başlığı.
  \item \textbf{Hiçbir metin atlanmaz:} her cümle, dipnot, tablo hücresi ve
        şekil altyazısı çevrilir.
  \item Bölüm ve ortam sayısı kaynakla \textbf{birebir aynı} kalır. Derleme
        adımı bölüm sayısını kaynakla karşılaştırır ve uyuşmazlıkta uyarır.
  \item Büyük kitaplar bölüm sınırlarından parçalanır (parça başına en çok
        yaklaşık 80 KB önerilir); tüm parçalar aynı terim listesini kullanır.
  \item Başlıklara ``Bölüm X'' yazılmaz; şablon bölüm numarasını kendisi ekler.
  \item Türkçe karakterler doğru yazılır (ç, ğ, ı, İ, ö, ş, ü); dosya
        \textbf{UTF-8} olarak kaydedilir.
\end{itemize}

\subsection{Adım 3 --- Derleyin}
\begin{enumerate}
  \item \dugme{2 · Çevrilmiş TeX'i derle} düğmesine basın. Bu düğme, klasörde
        \dosya{EXPORT\_TR.tex} görülene kadar soluk kalır. Dosya
        kaydedildikten sonra düğme birkaç saniye içinde kendiliğinden açılır;
        açılmazsa sol kenardaki \dugme{Kitap adı}'nın TeX'i çıkarırken
        kullandığınızla aynı olduğuna bakın.
  \item İş bitince \dugme{PDF'yi aç} (PDF'yi varsayılan uygulamayla açar) ve
        \dugme{İndir} düğmeleri çıkar. Türkçe PDF kitap klasörüne
        \dosya{<kitap\_adı>\_TR.pdf} olarak da kaydedilmiştir.
\end{enumerate}

Derleme \emph{manuel çeviri kipinde} çalışır: yapısal denetimler (ortam
dengesi, yer tutucu bütünlüğü, derleme) aynen yapılır, ama çevirmenin sözcük
tercihlerinin üzerine yazılmaz.

\subsection{Yapay zekâ asistanıyla (Claude Code)}\label{sec:claude}
TeX çıkarıldıktan sonra kart, çevrilecek dosyanın ve bu kitaba özel çeviri
talimatının (kitap klasöründeki \dosya{CEVIRI\_TALIMATI.md}) tam yollarını ve
``Claude Code'a şunu yaz:'' başlığı altında hazır bir istem gösterir. İstemi
Claude Code'a verin; asistan talimatı okur, metni bölümlere böler, ortak bir
terim sözlüğü çıkarır, çevirir ve \dosya{EXPORT\_TR.tex}'i kaydeder. Ardından
arayüzde \dugme{2 · Çevrilmiş TeX'i derle} düğmesine basın.

Kaynak koddan kurulumda proje klasöründe Claude Code'u açıp yalnızca
\textbf{``bu kitabı çevir''} demeniz de yeterlidir: Claude, depodaki
\dosya{CEVIRI\_TALIMATI.md} dosyasını okur ve akışı (çıkar, çevir, derle)
kendisi yürütür.

%% ===================== YEREL MODEL ======================================
\section{Yerel Model ile Çevir (Beta)}

\begin{enumerate}
  \item Orta alana kitabı (PDF ya da EPUB) yükleyin. Kitap bu yolun
        desteklediği kaynak dilde olmalıdır (Bölüm~\ref{sec:ikiyol}'deki
        tablo; dil denetimi: Bölüm~\ref{sec:dil}).
  \item \dugme{Yerel model ile çevir} kartında \dugme{Kaynak dil} listesine
        bakın: kitabın dili kendiliğinden ölçülür ve güveniyle birlikte yazar.
        Tespit belirsizse \dugme{Çevir}den önce dili seçmeniz istenir; tespit
        yanlışsa dili buradan değiştirin.
  \item \dugme{Model} listesinden çeviri modelini seçin (önerilen ve
        varsayılan: HY-MT 1.5 7B; listede ``önerilen'' diye işaretlidir). Liste, arayüz ilk açıldığında bir kez
        kendiliğinden taranır; bir modeli yeni hazırladıysanız ya da bir
        sunucuyu yeni açtıysanız \dugme{Modelleri tara} düğmesine basın
        (Bölüm~\ref{sec:modeller}). Kullanılamayan bir model seçiliyse
        \dugme{Çevir} düğmesi soluk kalır ve nedeni altında yazar.
  \item \dugme{Çevir} düğmesine basın. Model önce birkaç saniyelik son bir
        denetimden geçer; sunucu kapandıysa ya da model artık yoksa iş
        başlamaz ve nedeni yazılır. Denetim geçerse çıkarım, çeviri, kalite
        denetimi ve dizgi kendiliğinden ilerler; ilerleme panelinde çeviriyi
        yapan model yazar.
  \item İş bitince \dugme{PDF'yi aç} ve \dugme{İndir} düğmeleri çıkar. PDF
        kitap klasöründe \dosya{compiled/\allowbreak<kitap\_adı>\_TR.pdf}
        olarak kaydedilir.
\end{enumerate}

\begin{notkutu}
\textbf{Makine çevirisi künyesi:} Bu yolun PDF ve EPUB çıktısının başında
(kitabın özgün kapağı basıldıysa onun arkasında) kısa bir künye bulunur:
metnin bu bilgisayardaki hangi yerel modelle çevrildiğini, gözden geçirilmemiş
bir taslak olduğunu, kaynak dili ve çeviri tarihini yazar. ``PDF'den TeX
çıkar'' çıktısında künye yoktur. Makine çevirisini kamuya açık bir yerde
paylaşacaksanız künyeyi silmeyin ve kullandığınız modelin lisansındaki
koşullara bakın (lisans metinlerinin bağlantıları: Bölüm~\ref{sec:lisans}).
\end{notkutu}

Bu yolun çıktısı ön okuma ve taslak içindir (bkz.\ Bölüm~\ref{sec:ikiyol}'deki
dürüst not). Kalite denetiminin eşikleri HY-MT 1.5 7B'ye göre ayarlıdır; başka
bir modelle çevrilen ilk kitapta \dosya{qc\_report.txt} dosyasını gözden
geçirin.

\subsection{Dil denetimi}\label{sec:dil}
Yerel model kitabın dilini iki kapıda ölçer:
\begin{enumerate}
  \item \textbf{Ön kapı (çıkarımdan önce):} PDF'nin metin katmanına ya da
        EPUB metnine bakar. Metin katmanı olan ve dili bu yolda
        desteklenmeyen bir kitapta iş, uzun süren çıkarım başlamadan
        saniyeler içinde durur. Metin katmanı yoksa (taranmış kitap) ya da
        sonuç belirsizse ön kapı geçer ve kararı tam kapıya bırakır.
  \item \textbf{Tam kapı (çıkarımdan sonra):} Çıkarılan metni (taranmış
        kitapta OCR sonucunu) ölçer ve çeviri başlamadan karar verir.
        Ölçülen dil desteklenen dillerden biriyse çeviriye aktarılır: modele
        verilen talimat o dilden Türkçeye çevirmesini söyler, çevrilmeden
        kalan satırlar o dilde aranır. Tespit belirsizse (ölçülebilir
        düzyazı az ya da güven düşük) iş durmaz, kaynak İngilizce varsayılır
        ve \dosya{pipeline.log} dosyasına uyarı yazılır.
\end{enumerate}
Kartta dili elle seçtiyseniz (tespitten farklı bir dil ya da tespit
belirsizken yapılan seçim) iki kapı da atlanır ve çeviri seçtiğiniz dille
yapılır.

Kitabın desteklenmeyen bir dilde (ör.\ Çince ya da Portekizce) olduğu
güvenle ölçülürse iki kapı da işi \textbf{çıkış kodu 4} ile durdurur; çeviri
hiç başlamaz. Arayüz bu durumda ``Yerel model bu kitabı çevirmedi: kitabın
dili desteklenmiyor (kaynak dil kapısı, çıkış kodu 4).'' iletisini, altında
tespit edilen dili ve desteklenen dillerin listesini gösterir.

\begin{notkutu}
\textbf{Tespit yanlışsa:} kitabın dili aslında destekleniyorsa kartın
\dugme{Kaynak dil} listesinden doğru dili seçip yeniden çevirin. Komut
satırında (bkz.\ Bölüm~\ref{sec:komut}) \dosya{-{}-kaynak-dil} bayrağı dili
elle verir (ör.\ \dosya{-{}-kaynak-dil de}), \dosya{-{}-dil-kontrolu-yok} ise
kaynağı İngilizce varsayar. İki bayrak da iki kapıyı birden atlar;
\dosya{-{}-dil-kontrolu-yok} arayüzde yoktur ve yalnızca tespit yanlış
olduğunda kullanılır.
\end{notkutu}

%% ===================== MODEL KAYNAKLARI =================================
\section{Model Kaynakları}\label{sec:modeller}

Önerilen ve varsayılan çeviri modeli \textbf{HY-MT 1.5 7B}'dir (MLX, kimliği
\dosya{hymt-7b}): en çok sınanan model budur ve kalite denetiminin eşikleri
onun çıktısına göre ayarlıdır. DMG kurulumunda Bileşenler ekranından, kaynak
koddan kurulumda Bölüm~\ref{sec:kaynak}'teki komutla kurulur. Onun yerine dört
kaynaktan bir model seçilebilir; bu kaynaklar desteklenir ama \textbf{beta
düzeyindedir ve daha az sınanmıştır}. Her kaynakta akış aynıdır:
\textbf{hazırla} \(\rightarrow\) \dugme{Modelleri tara} \(\rightarrow\)
\textbf{seç} \(\rightarrow\) \dugme{Çevir}.

\vspace{6pt}
\renewcommand{\arraystretch}{1.35}
\noindent\begin{tabular}{@{}L{2.6cm}L{7.6cm}L{5.0cm}@{}}
\toprule
\textbf{Kaynak} & \textbf{Hazırlık} & \textbf{Kimlik örneği} \\
\midrule
MLX klasörü & MLX biçimli model \dosya{models/} altına konur
(\dosya{config.json}, \dosya{*.safetensors} ve tokenizer dosyalarıyla) &
\dosya{mlx:qwen3-8b-mlx} \\
Ollama & Ollama kurulur ve model indirilir (\dosya{ollama pull qwen3:8b});
uygulama açılır ya da \dosya{ollama serve} çalıştırılır &
\dosya{ollama:qwen3:8b} \\
LM Studio & Model LM Studio'da indirilip yüklenir; Developer bölümünden yerel
sunucu başlatılır ya da \dosya{lms server start} çalıştırılır &
\dosya{lmstudio:qwen/\allowbreak qwen3-8b} \\
GGUF dosyası & \dosya{.gguf} dosyası \dosya{models/} altına konur; doğrudan
açmak için llama-cpp-python gerekir. Dosya Ollama ya da LM Studio'ya da
yüklenebilir & \dosya{gguf:Qwen3-8B-\allowbreak Q4\_K\_M.gguf} \\
\bottomrule
\end{tabular}
\vspace{8pt}

Kitapçevir Ollama'yı ve LM Studio'yu bu programların varsayılan
adreslerinde arar; programların portunu değiştirmediyseniz hiçbir ayar
gerekmez. Adresler tek yerde tanımlıdır: \dosya{config/\allowbreak
defaults.py} dosyasındaki \dosya{OLLAMA\_\allowbreak DEFAULT\_\allowbreak URL}
ve \dosya{LMSTUDIO\_\allowbreak DEFAULT\_\allowbreak URL}. llama.cpp sunucusu, Jan,
vLLM gibi OpenAI uyumlu başka bir yerel sunucu da kullanılabilir: adresi
\dosya{config/\allowbreak mt\_sources.json} dosyasına yazılır, kimliği
\dosya{openai:<adres>|<model>} biçimindedir (ör.\
\dosya{openai:http://127.0.0.1:8080/v1|qwen}).

\subsection{Arayüzde seçim}
\begin{itemize}
  \item Listede yalnızca taramada gerçekten bulunan modeller görünür.
        Sunucusu kapalı bir kaynaktan model listelenmez; kartta o kaynağın
        durumu yazar (ör.\ ``Ollama sunucusu kapalı ya da erişilemiyor
        (\ldots)''). Yerleşik HY-MT 1.5 7B listede her zaman bulunur.
  \item Her seçenekte modelin adı ve kaynağı, düşünme durumu, biliniyorsa
        ayrıntısı (parametre sayısı, niceleme, mimari) ve boyutu yazar; HY-MT
        1.5 7B ``önerilen'' diye işaretlidir. Kullanılamayan
        model ``kullanılamaz'' diye işaretlidir; seçilirse \dugme{Çevir}
        soluk kalır ve nedeni altında yazar.
  \item Seçim hatırlanır. Son seçilen model bir sonraki taramada
        bulunamazsa varsayılan model seçilir ve bu durum kartta yazar. İş
        sürerken model listesi ve \dugme{Modelleri tara} kapalıdır.
  \item Model çeviri sırasında kullanılamaz hâle gelirse (ör.\ sunucu
        kapandı) iş \textbf{çıkış kodu 5} ile biter ve arayüz ``Seçili model
        kullanılamadı; nedeni aşağıda.'' iletisini gösterir.
\end{itemize}

\subsection{Başka klasördeki model: sembolik bağ}
Kitapçevir yalnızca \dosya{models/} klasörünü tarar (DMG kurulumunda
\dosya{veri/models/}); LM Studio'nun ya da Hugging Face'in indirme klasörleri
taranmaz. Başka yerdeki bir modeli kopyalamadan kullanmak için
\dosya{models/} altına sembolik bağ koyun (kaynak koddan kurulumda proje
klasöründe):

\begin{komut}
ln -s \textasciitilde/.lmstudio/models/<yayıncı>/<Model-MLX> models/<Model-MLX>
\end{komut}

DMG kurulumunda bağın hedefi
\dosya{\textasciitilde/Library/\allowbreak Application Support/\allowbreak Kitapçevir/\allowbreak veri/\allowbreak models/<Model-MLX>}
olur.

\subsection{GGUF dosyasını doğrudan açmak: llama-cpp-python}
llama-cpp-python \textbf{isteğe bağlı} bir bağımlılıktır. DMG kurulumunda
Bileşenler ekranındaki \textbf{GGUF desteği} onu kurar. Kaynak koddan
kurulumda \dosya{requirements.txt} onu kurmaz. Kurulu değilse uygulama yine
açılır, GGUF dosyaları listede ``kullanılamaz'' görünür ve kurulum komutu
yazılır. Kaynak koddan kurulumda Apple Silicon'da Metal desteğiyle kurulum
(proje klasöründe):

\begin{komut}
CMAKE\_ARGS="-DGGML\_METAL=on" ./venv/bin/pip install llama-cpp-python
\end{komut}

\subsection{Düşünen modeller}
Qwen3, gpt-oss gibi modeller yanıttan önce ``düşünür''. Kural şudur:
\textbf{düşünme resmî yolla kapatılabiliyorsa kapatılır; kapatılamıyorsa model
düşünerek çevirir.} HY-MT 1.5 7B düşünmeyen bir modeldir.
\begin{itemize}
  \item \textbf{Kapatılabiliyorsa:} model düşünmeden çevirir (ör.\ Ollama'da
        isteğe düşünmeyi kapatan alan eklenir, LM Studio'da düşünme kapalı
        istenir, MLX modelinde sohbet şablonuna kapatma değişkeni geçer).
  \item \textbf{Kapatılamıyorsa:} model düşünerek çevirir; her parçaya ayrı
        bir düşünme payı ayrılır ve iş daha yavaş sürer. Pay dolup cevap
        gelmeyen parça bir kez yeniden denenir; yine cevap yoksa o parça
        özgün dilinde kalır ve işaretlenir (kalite raporunda görünür), çeviri
        sürer.
  \item \textbf{Cevap üretemiyorsa:} model parçaların çoğunda cevap
        üretemiyorsa çeviri ``Bu model düşünmesini bitirip çeviri
        üretemiyor: \ldots'' iletisiyle, \textbf{çıkış kodu 5} ile durur ve
        önerilen model adıyla başka bir model seçmeniz istenir.
\end{itemize}
Listede her modelin düşünme durumu yazar (``düşünme yok'', ``düşünme kapalı
çalışır'', ``düşünen model --- düşünerek çevirir, daha yavaştır'' ya da
``düşünme durumu bilinmiyor''); kart seçili modelinkini de gösterir. Düşünme
durumu bilinmeyen ama düşündüğü bilinen model için kartta ``düşünen model;
yavaş olabilir'' notu çıkar.
Ayrıntı ve ayarlar: \dosya{docs/\allowbreak YENI\_MODEL\_EKLEME.md}, Bölüm~6.

\subsection{Bilinmesi gerekenler}
\begin{itemize}
  \item \textbf{Yalnızca bu bilgisayar:} sunuculara yalnızca yerel adresle
        (\dosya{127.0.0.1}, \dosya{localhost}, \dosya{::1}) bağlanılır; API
        anahtarı desteği yoktur.
  \item \textbf{Dil:} hedef dil yalnızca Türkçedir. Kaynak dil
        Bölüm~\ref{sec:ikiyol}'deki tabloda yazar; dil kapıları her modelde
        aynen çalışır.
  \item \textbf{Çeviri belleği model başına ayrıdır:} her model yalnızca
        kendi önceki çevirilerini kullanır; HY-MT'nin mevcut kayıtları olduğu
        gibi korunur.
  \item \textbf{Bellek:} Ollama / LM Studio'da yüklü modeller, uygulamanın
        kendi sürecindeki model ve PDF çıkarımı aynı birleşik belleği
        paylaşır; kullanmadığınız modelleri sunucuda boşaltın.
  \item \textbf{Ayar dosyaları} (isteğe bağlı, git dışı): çoğu zaman
        gerekmez. Sunucu varsayılandan farklı bir portta çalışıyorsa ya da
        llama.cpp sunucusu gibi ek bir sunucu kullanılacaksa adresi
        \dosya{config/\allowbreak mt\_sources.json} dosyasına; modele kısa
        bir ad ya da bağlam uzunluğu gibi bir ayar gerekiyorsa kaydı
        \dosya{config/\allowbreak mt\_engines.json} dosyasına elle yazılır.
        Örnek dosya yoktur; kısa örnekler
        \dosya{docs/\allowbreak YENI\_MODEL\_EKLEME.md} dosyasındadır.
        Tarama klasörü bu dosyalarla eklenemez.
\end{itemize}

Kimlik biçimlerinin, ayar alanlarının ve bütün hata iletilerinin ayrıntılı
anlatımı: \dosya{docs/\allowbreak YENI\_MODEL\_EKLEME.md}.

%% ===================== İLERLEME =========================================
\section{İlerleme, İptal ve Sonuç}

\begin{itemize}
  \item \textbf{İlerleme paneli} işin adını (``TeX çıkarılıyor'', ``Çeviri
        sürüyor'', ``Kitap diziliyor''), o anki fazı ve geçen süreyi
        gösterir. İlerleme çubuğu yalnızca pipeline gerçek bir sayaç
        verdiğinde (segment ya da sayfa sayısı) çizilir; sayacı olmayan
        aşamada çubuk yoktur.
  \item \textbf{İptal:} Paneldeki \dugme{İptal} düğmesi işi ve açtığı alt
        süreçleri durdurur.
  \item \textbf{Başarılı sonuç:} ``PDF hazır'' iletisi, \dugme{PDF'yi aç} ve
        \dugme{İndir} düğmeleri ve PDF'nin dosya yolu.
  \item \textbf{Kalite eşiği altında:} ``PDF üretildi ama kalite denetimi
        eşiğin altında kaldı'' iletisi çıkar. PDF yine açılabilir; ayrıntı
        kitap klasöründeki \dosya{qc\_report.txt} dosyasındadır.
  \item \textbf{Hata:} Hata iletisi ve ``Son çıktı'' bölümünde işin son
        satırları gösterilir. Tam kayıt kitap klasöründeki
        \dosya{pipeline.log} dosyasındadır.
\end{itemize}

%% ===================== ARAÇLAR ==========================================
\section{Araçlar}\label{sec:araclar}

\begin{description}[style=nextline]
  \item[Terimler] Seçili kitabın kendi sözlüğünü (kitap klasöründeki
        \dosya{glossary\_shared.tsv}; bkz.\ Bölüm~\ref{sec:sozluk}) gösterir.
        Kaynak dildeki terimde ya da Türkçe karşılıkta arar; tam eşleşme yoksa
        en yakın terimi benzerlik oranıyla önerir. Kitap seçilmemişse ya da
        kitabın sözlüğü henüz yoksa sekme bunu söyler.
  \item[Bellek] Çeviri belleğindeki kayıt ve onaylı kayıt sayısını gösterir;
        yazdığınız metne benzeyen kayıtları (en çok beş) kaynak ve
        çevirisiyle listeler. Bu sekme belleğe yalnızca bakar, bir şey
        değiştirmez.
  \item[Karşılaştırma] Orijinal ve çevrilmiş PDF'yi yükleyip
        \dugme{Karşılaştır}'a basınca iki PDF'yi sayfa sayfa piksel düzeyinde
        karşılaştırır: karşılaştırılan sayfa sayısı, ortalama fark ve sorunlu
        sayfalar gösterilir, HTML rapor indirilebilir. ``Karşılaştırma
        ayarları''ndan sayfa sınırı (varsayılan 50 ya da tüm sayfalar), DPI
        (72--300, varsayılan 150) ve görsellerin rapora eklenip eklenmeyeceği
        seçilir.
\end{description}

\subsection{Kitap sözlüğü}\label{sec:sozluk}
Terim tutarlılığı her kitabın kendi sözlüğüyle sağlanır; bütün kitaplara
uygulanan sabit bir genel alan sözlüğü yoktur. Sözlük kitap klasöründeki
\dosya{glossary\_shared.tsv} dosyasıdır: UTF-8, her satır ``kaynak terim,
sekme, Türkçe karşılık'' (isteğe bağlı üçüncü sütun not); kaynak terim
kitabın kendi dilindedir. İki yol aynı biçimi kullanır.

\begin{itemize}
  \item \textbf{PDF'den TeX çıkar:} Sözlüğü çeviriyi yapan çıkarır. Yapay
        zekâ asistanı, kitaba özel \dosya{CEVIRI\_TALIMATI.md} gereği
        çeviriden önce kitabın tekrar eden teknik terimlerini (50--100 terim)
        bu dosyaya yazar ve bütün parçalarda kullanır.
  \item \textbf{Yerel model ile çevir:} Kitap klasöründe sözlük yoksa
        uygulama, çeviri bittikten sonra kitabın sık geçen terimlerini kaynak
        dilde (altı dilin hepsinde) çıkarır, her terimi bir kez, geçtiği
        cümleyle birlikte seçili modele çevirtir ve dosyaya yazar. Kaynak dili
        İngilizce olan kitapta, elle doğrulanmış başlangıç terimlerinden
        (\dosya{glossary/\allowbreak curated.json}) kitapta geçenler önce
        gelir. Ardından çeviri bu sözlükle denetlenir: çıktıda özgün dilinde
        kalmış terim sözlükteki karşılığıyla değiştirilir; aynı terimin
        farklı çevrildiği yerler kalite raporunda düşük önemli uyarı olarak
        görünür.
  \item \textbf{Var olan sözlük korunur:} Dosya zaten varsa (sizin, TeX
        yolunun ya da önceki bir çevirinin yazdığı) üzerine yazılmaz, olduğu
        gibi kullanılır.
\end{itemize}

%% ===================== KOMUT SATIRI =====================================
\section{Komut Satırından Kullanım}\label{sec:komut}

Arayüz, işleri aşağıdaki betikleri ayrı bir süreçte çalıştırarak yapar;
kaynak koddan kurulumda aynı işler doğrudan terminalden de yapılabilir. Sanal
ortam etkin olmalıdır (\dosya{source venv/bin/activate}).

\subsection{PDF'den TeX çıkar}
\begin{komut}
python 1\_cikar.py kitap.pdf \ \ \ \ \ \ \ \ \ \ \ \ \ \ \ \ \ \ \# \(\rightarrow\) output/<ad>/EXPORT.tex\\
\# EXPORT.tex'i çevirin \(\rightarrow\) aynı klasöre EXPORT\_TR.tex olarak kaydedin\\
python 2\_derle.py output/<ad>/EXPORT\_TR.tex \ \ \# \(\rightarrow\) Türkçe PDF
\end{komut}

\dosya{1\_cikar.py} aynı adlı çıktı klasöründe farklı bir kitap bulursa durur
ve \dosya{-{}-name} ile farklı bir ad vermenizi ister. \dosya{2\_derle.py}
için \dosya{-{}-output-dir output/<ad>} seçeneği, kimlik satırları silinmiş
ya da dosya taşınmışsa kitabın klasörünü elle gösterir.

\subsection{Yerel model ile çevir}
\begin{komut}
python translate\_book.py kitap.pdf \ \ \ \ \# PDF \(\rightarrow\) Türkçe PDF\\
python translate\_book.py kitap.epub \ \ \ \# EPUB \(\rightarrow\) Türkçe PDF (+ Türkçe EPUB)
\end{komut}

Sık kullanılan seçenekler:
\begin{itemize}
  \item \dosya{-{}-name "Kitap Adı"}: çıktı klasörünün ve PDF'nin adı.
  \item \dosya{-{}-resume}: çökme sonrasında kaldığı yerden devam eder.
  \item \dosya{-{}-book-type}: kitap türü; türe göre kuralları (ör.\ kalite
        denetimi) etkiler.
  \item \dosya{-{}-skip-compile}: yalnızca \dosya{.tex} üretir, PDF derlemez.
  \item \dosya{-{}-batch}: birden çok kitabı sırayla işler.
  \item \dosya{-{}-mt-engine KİMLİK}: çeviri modelini seçer (varsayılan
        \dosya{hymt-7b}); ör.\ \dosya{ollama:qwen3:8b}. Biçimler:
        Bölüm~\ref{sec:modeller}. Arayüzün son seçimini okumaz.
  \item \dosya{-{}-dil-kontrolu-yok}: kaynak dilin iki kapısını birden
        (çıkarımdan önceki ön kapı ve çıkarımdan sonraki tam kapı) atlar ve
        kaynağı İngilizce varsayar. Yalnızca dil tespiti yanlışsa, yani
        kitabın dili aslında destekleniyorken kullanılır; arayüzde yoktur
        (bkz.\ Bölüm~\ref{sec:dil}).
  \item \dosya{-{}-kaynak-dil KOD}: kitabın kaynak dilini elle verir
        (\dosya{en}, \dosya{de}, \dosya{ru}, \dosya{fr}, \dosya{es} ya da
        \dosya{it}); iki dil kapısını da atlar. Arayüzdeki karşılığı kartın
        \dugme{Kaynak dil} listesidir.
\end{itemize}
Kitabın desteklenmeyen bir dilde olduğu güvenle ölçülürse
\dosya{translate\_book.py} çeviriye başlamadan \textbf{çıkış kodu 4} ile
durur.
Çıktı \dosya{output/<ad>/} altında toplanır: \dosya{source/} (ara dosyalar),
\dosya{compiled/} (PDF), \dosya{qc\_report.txt} (kalite raporu),
\dosya{pipeline.log} (iş kaydı).

Seçili model kullanılamıyorsa (kimlik geçersiz, sunucu kapalı, model yok,
bu Mac'in belleği yetmiyor, çeviri sırasında sunucu bağlantısı koptu ya da
düşünen model parçaların çoğunda cevap üretemedi) \dosya{translate\_book.py}
\textbf{çıkış kodu 5} ile durur; model ve bellek çıkarımdan önce denetlendiği
için uzun PDF çıkarımı boşuna beklenmez (bellek: Bölüm~\ref{sec:donanim}).
Model seçerek çeviri ve bulunan modellerin listesi:

\begin{komut}
python translate\_book.py kitap.pdf --mt-engine ollama:qwen3:8b\\
python -m core.mt\_engines.discovery \ \ \ \ \# bulunan modeller, sunucu durumları
\end{komut}

%% ===================== SORUN GİDERME ====================================
\section{Sık Karşılaşılan Sorunlar}

% Tablo uzadıkça sayfaya sığmayabilir: longtable sayfa sonunda bölünür,
% başlık satırı yeni sayfada tekrarlanır.
\renewcommand{\arraystretch}{1.35}
\setlength{\LTleft}{0pt}
\setlength{\LTright}{\fill}
\begin{longtable}{@{}L{5.7cm}L{9.7cm}@{}}
\toprule
\textbf{Belirti} & \textbf{Çözüm} \\
\midrule
\endhead
macOS uygulamayı açmıyor (tanınmayan geliştirici uyarısı) & Uygulama
imzasızdır; Bölüm~\ref{sec:ilkacilis}'teki adımları izleyin: Finder'da sağ tık
\ok{} Aç; macOS 15 ve sonrasında Sistem Ayarları \ok{} Gizlilik ve Güvenlik
\ok{} Yine de Aç. \\
Açılış ekranında ``Kitapçevir'i Uygulamalar klasörüne taşıyıp oradan açman
önerilir.'' & Uygulama DMG'nin içinden ya da taşınmadan açıldı. Uygulamayı
Uygulamalar klasörüne sürükleyip oradan açın. \\
Sihirbaz: ``Datalab model lisansı kabul edilmeden kurulum sürdürülemez'' &
PDF desteği zorunlu bileşendir ve modelleri olmadan çalışmaz. Kurulumu
sürdürmek için lisansı okuyup kabul edin. \\
Sihirbaz: ``Boş disk: \ldots{} Yetmiyor: \ldots{} daha gerekiyor.'' & Yer
açın ya da isteğe bağlı bileşenleri (ör.\ HY-MT) seçmeden kurun; onları sonra
Bileşenler ekranından ekleyebilirsiniz. \\
Sihirbaz: ``Bağlantı yok.'' & İnternet bağlantısını açıp \dugme{Yeniden
dene}'ye basın. Yarım inen dosyalar saklanır. \\
Sihirbazda HY-MT seçilemiyor & Bu Mac'in belleği HY-MT'ye yetmiyor
(Bölüm~\ref{sec:donanim}). ``PDF'den TeX çıkar'' yolunu kullanın. \\
``Kurulum başlatılamadı'' ekranı & Ekrandaki iletiyi ve günlüğü okuyun
(\dugme{Göster} günlük dosyasını açar). Günlükler
\dosya{\textasciitilde/Library/\allowbreak Logs/\allowbreak Kitapçevir/}
altındadır. \\
Arayüz ``yalnızca macOS'ta çalışır'' ya da ``Apple Silicon (M serisi) bir
Mac ister'' yazıp duruyor & Kitapçevir yalnızca Apple Silicon Mac'te çalışır.
Kaynak koddan kurulumda Apple Silicon Mac'te bile bu ileti çıkıyorsa Python
Rosetta altında (x86\_64) çalışıyordur; sanal ortamı arm64 Python ile yeniden
kurun. \\
\dugme{Çevir} soluk, seçili model ``kullanılamaz'' & Nedeni kartta yazar.
HY-MT kurulu değilse DMG kurulumunda Bileşenler ekranından indirin; kaynak
koddan kurulumda \dosya{mlx\_lm.convert} komutunu çalıştırın
(Bölüm~\ref{sec:kaynak}); diğer modeller için Bölüm~\ref{sec:modeller}. \\
Ollama ya da LM Studio modeli listede yok & Sunucu kapalıdır; kartta durumu
yazar. Ollama uygulamasını açın (ya da \dosya{ollama serve}) veya LM
Studio'da yerel sunucuyu başlatın (\dosya{lms server start}), sonra
\dugme{Modelleri tara}. LM Studio'da model yüklü değilse ve isteğe bağlı
yükleme (Just-in-Time) kapalıysa model listelenmez. Sunucu varsayılandan
farklı bir portta çalışıyorsa adresi \dosya{config/\allowbreak
mt\_sources.json} dosyasına yazın. \\
Başka klasördeki bir model listede yok & Yalnızca \dosya{models/} taranır;
\dosya{models/} altına sembolik bağ koyun (Bölüm~\ref{sec:modeller}). \\
GGUF dosyası ``kullanılamaz'', kurulum komutu yazıyor & llama-cpp-python
kurulu değil (isteğe bağlıdır). DMG kurulumunda Bileşenler ekranından GGUF
desteğini indirin; kaynak koddan kurulumda komutla kurun ya da dosyayı
Ollama / LM Studio üzerinden kullanın. \\
Çıkış kodu 5, ``Seçili model kullanılamadı; nedeni aşağıda.'' & Model
çeviriden önce ya da çeviri sırasında kullanılamaz hâle geldi (sunucu kapalı,
model yok, bağlantı koptu). Nedeni iletinin altında ve \dosya{pipeline.log}
dosyasında yazar. \\
``Bu model düşünmesini bitirip çeviri üretemiyor: \ldots'' (çıkış kodu 5) &
Düşünmesi kapatılamayan model parçaların çoğunda cevap üretemedi. İletide
önerilen modeli (HY-MT 1.5 7B) ya da düşünmeyen başka bir model seçin. \\
``Yalnız bu bilgisayardaki sunuculara bağlanılır; \ldots'' & Başka bir
bilgisayardaki sunucu desteklenmez; adres \dosya{127.0.0.1},
\dosya{localhost} ya da \dosya{::1} olmalı. \\
\dugme{2 · Çevrilmiş TeX'i derle} soluk & Çeviri, \dosya{EXPORT.tex} ile aynı
klasörde tam olarak \dosya{EXPORT\_TR.tex} adıyla durmalıdır. Dosya
kaydedildikten sonra düğme birkaç saniye içinde kendiliğinden açılır;
açılmazsa sol kenardaki \dugme{Kitap adı}'nın TeX'i çıkarırken
kullandığınızla aynı olduğuna bakın. \\
Kitap yüklenmiyor, boyut sınırı iletisi çıkıyor & Yükleme sınırı 500 MB'tır;
daha büyük kitaplar için kaynak koddan kurulumda komut satırını kullanın. \\
\dosya{xelatex: command not found} (kaynak koddan kurulum) & MacTeX'i kurun
ve yeni bir terminal açın. \\
Derleme: ``hangi kitaba ait olduğu çözülemedi'' & Çevrilmiş dosyada
\dosya{\% KITAPCEVIR-META} satırları silinmiş olabilir. \dosya{2\_derle.py}'ye
\dosya{-{}-output-dir output/<ad>} verin. \\
Derleme: ``full\_book.md bulunamadı'' & Önce TeX'i çıkarın (aynı kitap
adıyla). \\
Derleme: ``UTF-8 değil'' & Dosyayı UTF-8 olarak kaydedin. \\
``FARKLI bir kitap var'' & Aynı adlı çıktı klasöründe başka bir kitap var;
sol kenarda (ya da \dosya{-{}-name} ile) farklı bir kitap adı verin. \\
``Yerel model bu kitabı çevirmedi: kitabın dili desteklenmiyor'' (çıkış
kodu 4) & Kitabın dili desteklenen dillerden biri değilse bu beklenen
davranıştır. Dil aslında destekleniyorsa tespit yanlıştır; kartın
\dugme{Kaynak dil} listesinden doğru dili seçip yeniden çevirin (komut
satırında \dosya{-{}-kaynak-dil}, kitabın dilinin koduyla; bkz.\
Bölüm~\ref{sec:dil}). \\
Yerel model, TeX yoluyla derlenmiş bir kitabı çevirmeyi reddediyor & Bu
bilinçli bir korumadır: yerel model önceki çevirinizin üzerine yazmaz.
Kitabı model ile de çevirmek istiyorsanız farklı bir kitap adı verin. \\
İlk PDF çok yavaş (kaynak koddan kurulum) & Marker modelleri ilk seferde
iniyor (yaklaşık 3,5 GB); sonraki kitaplarda bu indirme yapılmaz. DMG
kurulumunda modeller kurulum sırasında iner. \\
``Bu Mac yerel model için yetersiz: \ldots'' (yerel model kartı kapalı) & Bu
Mac 4 bit bir 7B modeli bile rahat çalıştıramaz
(Bölüm~\ref{sec:donanim}). Model gerektirmeyen ``PDF'den TeX çıkar'' yolunu
kullanın. \\
Model listede ``bu Mac'e sığmaz (gerekli \textasciitilde\ldots{} GB,
kullanılabilir \ldots{} GB)'' & Modelin tahmini gereksinimi bu Mac'in
kullanılabilir belleğini aşıyor. Daha küçük ya da daha çok nicemlenmiş
(ör.\ 4 bit) bir model seçin (Bölüm~\ref{sec:modeller}). \\
\bottomrule
\end{longtable}

%% ===================== LİSANSLAR ========================================
\section{Lisanslar}\label{sec:lisans}
Kitapçevir GNU GPL v3.0 ile dağıtılır; tam metin deponun kökündeki
\dosya{LICENSE} dosyasındadır. Kurulumda indirilen ya da kullanılan başlıca
üçüncü taraf bileşenler kendi lisanslarıyla gelir. Koşullar burada
özetlenmez; bağlayıcı olan, aşağıda bağlantısı verilen resmî lisans
metinleridir.

\begin{itemize}
  \item \textbf{Marker ve Surya} (PDF çıkarımı; marker-pdf 1.10.2, surya-ocr
        0.17.1): kod GPL-3.0-or-later; model ağırlıkları Datalab'ın ``AI Pubs
        Open RAIL-M (değiştirilmiş)'' lisansıyla sunulur. Metinler:
        \url{https://github.com/datalab-to/marker} (\dosya{LICENSE},
        \dosya{MODEL\_LICENSE}) ve \url{https://github.com/datalab-to/surya}.
        Ağırlıklar kurulumda Datalab'ın sunucusundan iner; sihirbaz
        indirmeden önce lisansı gösterip onay ister.
  \item \textbf{PyMuPDF} (PDF okuma ve sayfa işlemleri; GNU AGPL-3.0 ya da
        Artifex ticari lisansı) ve \textbf{EbookLib} (EPUB; GNU
        AGPL-3.0-or-later). Kaynak kodları ve lisans metinleri:
        \url{https://github.com/pymupdf/PyMuPDF} ve
        \url{https://github.com/aerkalov/ebooklib}.
  \item \textbf{HY-MT 1.5 7B} (isteğe bağlı çeviri modeli): Tencent HY
        Community License; metin:
        \url{https://huggingface.co/tencent/HY-MT1.5-7B}
        (\dosya{License.txt}). Kitapçevir, Berke Tezgöçen tarafından
        geliştirilen bağımsız bir projedir; Tencent ile bağlantılı değildir,
        Tencent tarafından desteklenmez ya da onaylanmaz.
  \item \textbf{DejaVu yazı tipleri} (\dosya{fonts/}): lisans metni
        \dosya{fonts/\allowbreak LICENSE-DejaVu.txt} dosyasındadır.
  \item \textbf{TeX:} tectonic (MIT;
        \url{https://github.com/tectonic-typesetting/tectonic}); tectonic'in
        indirdiği TeX paketleri ve MacTeX, TeX Live bileşenlerinin kendi
        özgür lisanslarıyla gelir.
\end{itemize}

Bileşenler ekranındaki \dugme{Lisanslar} bağlantısı uygulamanın lisans
klasörünü açar: Kitapçevir'in GPL metni, kaynak kodu bağlantıları, Datalab
ve Tencent lisanslarının resmî metinleri, DMG'deki Python'un
(python-build-standalone; içine derlenen OpenSSL, SQLite, Tcl/Tk gibi
kitaplıklar dahil) ve kurulumda inen uv, tectonic ve GGUF bileşenlerinin
kendi lisans dosyaları; \dosya{00\_OKU.txt} hepsini sıralar. Kurulan Python
paketlerinin lisansları kendi \dosya{*.dist-info} klasörlerindedir.

%% ===================== TELİF ============================================
\section{Telif Notu}
Bu sistem kişisel okuma için tasarlanmıştır. Üretilen çeviriler yalnızca
kişisel kullanım içindir (FSEK m.\ 38 --- kişisel kullanım için çoğaltma
istisnası); satılmaz veya dağıtılmaz. Kitap dosyaları ve çeviri çıktıları
depoya eklenmez.

\end{document}
